% !TEX root = ../mainfile.tex
% !TEX encoding = UTF-8 Unicode
% !TEX spellcheck = en_US

\chapter{Background: The Events of 2022--2026}
\label{ch:background}

This chapter narrates the episodes the empirical chapters analyse and fixes the dates used
throughout. Three are central---the 2025 U.S.\ tariff shock, the 2026 Iran/Hormuz crisis,
and the 2022 Russian invasion of Ukraine as an external control---together with two
within-crisis sub-events (the April~2025 tariff pause and the June~2026 ceasefire) that
serve as natural reversals.

\section{Liberation Day and the U.S.\ Tariff Shock}
\label{sec:bg_tariff}

On 2~April 2025, branded ``Liberation Day,'' the United States announced broad tariffs on
nearly all trading partners. The market reaction was severe and, for the dollar,
anomalous. Equities fell sharply, implied volatility spiked, and crude oil dropped as
markets priced weaker global demand---yet the dollar \emph{depreciated} rather than
attracting the usual flight-to-safety bid, sliding several percent against the euro, yen and
Swiss franc. Long-term Treasury yields, after an initial decline, backed up even as the
currency weakened, a co-movement widely read as evidence that the episode dented the
dollar's safe-haven premium rather than reinforcing it \parencite{corsetti2025,ecb2025fsr}.
A partial ninety-day tariff pause on 9~April 2025 tempered the equity drawdown but did not
restore the dollar, suggesting markets treated the shock as a regime change in U.S.\
trade policy rather than a transient surprise.

\section{The Israel--Iran Escalation and the Strait of Hormuz}
\label{sec:bg_iran}

The Gulf episode unfolded in stages. A first, twelve-day confrontation in June~2025
(Israeli and U.S.\ strikes on Iranian facilities, an Iranian response, and a rapid
ceasefire) lifted crude only briefly and left currencies little changed. Escalation resumed
in February~2026 and turned acute on 4~March 2026, when Iran moved to close the Strait of
Hormuz---the conduit for roughly a fifth of seaborne oil---triggering attacks on shipping
and a sharp spike in Brent \parencite{fattouh2026,kilian2026}. The dollar strengthened on
impact, the classical safe-haven response and the mirror image of the tariff episode. A
further round of U.S.\ strikes followed on 28~May 2026, and a sixty-day ceasefire announced
on 11~June 2026, committing to reopen the Strait, reversed much of the move in both crude and
the dollar---providing a clean within-crisis test of whether the premium that the threat put
into the dollar came back out on de-escalation.

\section{The Ukraine Control}
\label{sec:bg_ukraine}

The Russian invasion of Ukraine on 24~February 2022 is the external benchmark: a
military/oil-supply shock not authored by the United States. It produced the textbook
pattern---a sustained dollar appreciation alongside a jump in energy prices---and so anchors
what a ``classical'' safe-haven response looks like when the United States is a bystander
rather than a protagonist.

\section{Event Timeline}
\label{sec:bg_timeline}

\begin{table}[htbp]
\centering
\begin{threeparttable}
\caption{Event timeline and shock classification}
\label{tab:timeline}
\small
\begin{tabular}{l l l l}
\toprule
Date & Event & Type & Oil \\
\midrule
2022-02-24 & Russia invades Ukraine        & External military/oil (control) & $\uparrow$ \\
2025-04-02 & Liberation Day tariffs         & U.S.\ trade policy             & $\downarrow$ \\
\quad 2025-04-09 & \quad 90-day tariff pause & \quad sub-event                & $\downarrow$ \\
2025-06-13 & Twelve-day war (Israel/U.S.--Iran) & U.S.-implicated military   & $\uparrow$ (mild) \\
2026-02-28 & Hormuz crisis begins           & U.S.-implicated oil supply     & $\uparrow\uparrow$ \\
\quad 2026-03-04 & \quad Strait of Hormuz closed & \quad sub-event            & $\uparrow\uparrow$ \\
2026-05-28 & U.S.\ strikes on Iran          & U.S.-implicated military       & $\uparrow$ \\
2026-06-11 & Sixty-day ceasefire            & de-escalation                  & $\downarrow$ \\
\bottomrule
\end{tabular}
\begin{tablenotes}[flushleft]\footnotesize
\item \textit{Notes.} ``Type'' classifies each shock by origin and nature; the final column
gives the qualitative crude-oil response. Sub-events are reversals nested within their
parent episode.
\end{tablenotes}
\end{threeparttable}
\end{table}

\section{Key Variables Over the Sample}
\label{sec:bg_descriptive}

Figure~\ref{fig:overview} plots the broad dollar index, the VIX, crude oil and the
geopolitical-risk index over 2019--2026, with the four headline events shaded. The visual
record previews the central result: the dollar rises into the Ukraine and Hormuz bands but
falls through the Liberation Day band, while oil and risk gauges move with the geopolitical
episodes. The formal analysis follows in Chapters~\ref{ch:event_study}
and~\ref{ch:regression}.

\begin{figure}[htbp]
  \centering
  \includegraphics[width=\linewidth]{fig1_overview}
  \caption[Key variables, 2019--2026]{Key variables over the sample with the four main events
  shaded (Ukraine invasion, Liberation Day, twelve-day war, Hormuz crisis). Sources: FRED,
  LSEG, and Caldara and Iacoviello's Geopolitical Risk index \parencite{caldara2022}.}
  \label{fig:overview}
\end{figure}
